\section{积分的计算}
\label{sec:ru-integration-methods}

微积分基本定理把定积分的计算化为寻找原函数。本节先由导数表整理基本函数的原函数，再说明换元与分部积分怎样利用已知的原函数处理新的被积函数。

\subsection{基本函数的积分}
\label{sec:ru-basic-antiderivatives}

把导数表倒过来读，就得到基本函数的原函数。在所注明的区间上，
\[
 \begin{aligned}
 &\int x^\alpha\,\mathrm dx=\frac{x^{\alpha+1}}{\alpha+1}+C\ (\alpha\neq-1),&&\int\frac{\mathrm dx}x=\log|x|+C\ (x\neq0),\\
 &\int e^x\,\mathrm dx=e^x+C,&&\int a^x\,\mathrm dx=\frac{a^x}{\log a}+C\ (a>0,\ a\neq1),\\
 &\int\sin x\,\mathrm dx=-\cos x+C,&&\int\cos x\,\mathrm dx=\sin x+C,\\
 &\int\frac{\mathrm dx}{1+x^2}=\arctan x+C,&&\int\frac{\mathrm dx}{\sqrt{1-x^2}}=\arcsin x+C\ (|x|<1).
 \end{aligned}
\]
幂函数的公式在$\alpha=-1$时分母为零，这一情形单独由对数给出；$\log|x|$在$x<0$时的导数为$\frac{1}{-x}\cdot(-1)=\frac1x$，所以公式在两侧都成立。$x^\alpha$的允许区间与第~\ref{sec:ru-elementary-derivatives}节相同。

记号$\int f(x)\,\mathrm dx=G(x)+C$称为\term{不定积分}（indefinite integral），它表示$f$在一个区间上的全部原函数组成的一族函数，$C$取遍实常数；定积分$\int_a^bf$则是一个数。

若定义域由互不相连的几个区间组成，每个区间上可以取不同的常数，例如$1/x$在$(-\infty,0)$与$(0,+\infty)$上的原函数可以各取一个常数。
线性性质把复杂的被积函数拆成基本函数，例如
\[
 \int_1^2\Bigl(3x^2-\frac2x\Bigr)\,\mathrm dx=\bigl[x^3-2\log x\bigr]_1^2=7-2\log2.
\]

求出原函数后，可以求导回代核对，再检查积分区间是否落在公式成立的范围内。若忽略这一步，会得到错误的结果：$\int_{-1}^1x^{-2}\,\mathrm dx$按$[-1/x]_{-1}^1$计算得$-2$，
可是被积函数为正，积分不可能为负。问题在于$x^{-2}$在$0$处无界，不满足Newton--Leibniz公式的条件；这个积分要按第~\ref{sec:ru-improper-integrals}节的反常积分处理，实际上是发散的。

\subsection{换元法}
\label{sec:ru-substitution}

链式法则给出复合函数的导数，倒过来就是换元积分。
\begin{theorem}[换元公式]
\label{thm:ru-substitution}
设$\varphi$在$[\alpha,\beta]$上连续可微，$f$在包含$\varphi([\alpha,\beta])$的区间上连续，则
\[
 \int_{\varphi(\alpha)}^{\varphi(\beta)}f(u)\,\mathrm du=\int_\alpha^\beta f(\varphi(t))\,\varphi'(t)\,\mathrm dt.
\]
\end{theorem}
\begin{proof}
设$F$是$f$的原函数。由链式法则，$(F\circ\varphi)'=f(\varphi)\varphi'$，右侧被积函数连续，所以由Newton--Leibniz公式右边等于$F(\varphi(\beta))-F(\varphi(\alpha))$，左边也等于这个数。
\end{proof}
公式中的微分因子与积分端点一同变化：$\mathrm du$换成$\varphi'(t)\,\mathrm dt$，反映输入区间长度的伸缩；端点换成$\varphi$的对应值。定理不要求$\varphi$单调，$\varphi'<0$时端点的先后颠倒，由有向积分的约定自动处理符号。
例如$\int_0^12xe^{x^2}\,\mathrm dx$的被积函数形如$f(\varphi(x))\varphi'(x)$，取$\varphi(x)=x^2$、$f(u)=e^u$，得$\int_0^1e^u\,\mathrm du=e-1$。

另一种用法是引入新变量化简被积函数，这时要选定反函数的分支。计算四分之一单位圆的面积$\int_0^1\sqrt{1-x^2}\,\mathrm dx$，令$x=\sin\theta$。
新变量的范围取$[0,\pi/2]$，它使$\sin$从$0$递增到$1$，端点对应$0\mapsto0$、$\pi/2\mapsto1$，方向不变。在这一范围内$\cos\theta\geq0$，所以$\sqrt{1-\sin^2\theta}=\cos\theta$，开方不带负号。于是
\[
 \int_0^1\sqrt{1-x^2}\,\mathrm dx=\int_0^{\pi/2}\cos^2\theta\,\mathrm d\theta=\Bigl[\frac\theta2+\frac{\sin2\theta}4\Bigr]_0^{\pi/2}=\frac\pi4,
\]
其中用到$\cos^2\theta=(1+\cos2\theta)/2$。若把范围取成$[\pi/2,\pi]$，$\cos\theta\leq0$，开方必须写成$-\cos\theta$，同时端点方向也颠倒，结果仍是$\pi/4$；忽略其中一项就会出错。
一般可测函数在更一般的变换下的换元，放在第~\ref{chap:measurable-functions}章。

\subsection{分部积分法}
\label{sec:ru-integration-by-parts}

乘积的导数$(uv)'=u'v+uv'$含有两项。对它积分并移项，就能把其中一项的积分换成另一项。
\begin{theorem}[分部积分公式]
\label{thm:ru-integration-by-parts}
设$u,v$在$[a,b]$上连续可微，则
\[
 \int_a^bu(x)v'(x)\,\mathrm dx=\bigl[u(x)v(x)\bigr]_a^b-\int_a^bu'(x)v(x)\,\mathrm dx.
\]
\end{theorem}
\begin{proof}
由乘积法则，$uv$在$[a,b]$上连续可微且$(uv)'=u'v+uv'$，右侧两项都连续。由Newton--Leibniz公式，$\int_a^b(u'v+uv')=u(b)v(b)-u(a)v(a)$，再由积分的线性移项即得。
\end{proof}

使用时把被积函数拆成$u$与$v'$两个因子：求导后变简单的因子取作$u$，容易找到原函数的因子取作$v'$。


计算$\int_0^1xe^x\,\mathrm dx$，取$u=x$、$v'=e^x$，则$u'=1$、$v=e^x$，于是结果为$[xe^x]_0^1-\int_0^1e^x\,\mathrm dx=e-(e-1)=1$；边界项$[xe^x]_0^1=e$不能遗漏。
只有一个因子时，可以把$1$当作另一个因子：$\int_1^e\log x\,\mathrm dx$取$u=\log x$、$v'=1$，得$[x\log x]_1^e-\int_1^e1\,\mathrm dx=e-(e-1)=1$。
若反过来分配，例如对$\int x\sin x\,\mathrm dx$取$u=\sin x$、$v'=x$，剩下的积分含$x^2\cos x$，反而更复杂。

\paragraph{Taylor公式的积分余项}
分部积分的一个重要应用，是把Taylor公式的余项写成积分，从而去掉Lagrange余项中位置未知的中间点。设$f$在包含$a$与$a+h$的区间上$n+1$阶连续可微，令$\phi(t)=f(a+th)$，$t\in[0,1]$，则$\phi^{(k)}(t)=h^kf^{(k)}(a+th)$。
从$\phi(1)-\phi(0)=\int_0^1\phi'(t)\,\mathrm dt$出发，对$k\geq1$，在$\int_0^1\frac{(1-t)^{k-1}}{(k-1)!}\phi^{(k)}(t)\,\mathrm dt$中取$u=\phi^{(k)}$、$v=-(1-t)^k/k!$，分部积分给出
\[
 \int_0^1\frac{(1-t)^{k-1}}{(k-1)!}\phi^{(k)}(t)\,\mathrm dt=\frac{\phi^{(k)}(0)}{k!}+\int_0^1\frac{(1-t)^k}{k!}\phi^{(k+1)}(t)\,\mathrm dt.
\]
从$k=1$用到$k=n$，每一步分出一项$f^{(k)}(a)h^k/k!$，便得到带积分余项的Taylor公式
\begin{equation}
 f(a+h)=\sum_{k=0}^n\frac{f^{(k)}(a)}{k!}h^k+\frac{h^{n+1}}{n!}\int_0^1(1-t)^nf^{(n+1)}(a+th)\,\mathrm dt.
 \label{eq:ru-taylor-integral-remainder}
\end{equation}
若$|f^{(n+1)}|\leq M$，由$\int_0^1(1-t)^n\,\mathrm dt=1/(n+1)$，余项不超过$M|h|^{n+1}/(n+1)!$，与第~\ref{sec:ru-taylor-remainders}节的Lagrange余项给出同样的界。
积分余项不含未知点，可以直接夹逼。取$f=e^x$、$a=0$、$h=1$，余项为$\frac1{n!}\int_0^1(1-t)^ne^t\,\mathrm dt$，由$1\leq e^t\leq e$，它介于$1/(n+1)!$与$e/(n+1)!$之间；
$n=3$时$e-8/3\approx0.0516$，确实落在$1/24$与$e/24$之间。第~\ref{sec:ru-numerical-quadrature}节估计求积公式的局部误差时，同样要用Taylor展开把误差归结为高阶导数的界。
